%% kapa: exam=1 round=24 year=2013 period=1 form=A subject=economics subject_ko=경제원론 seq=9 number=49
%% source: Q-Net 감정평가사 기출문제 (cst00309) artlSeq=5073802 file=extracted/1008/5073802/2013년 제24회 감정평가사 1차시험 기출문제/1교시(A형).hwp
%% conversion: kordoc 4.24.51 -> tools/gen1.py (auto); review=proofread-sonnet; flags=none
\begin{question}{49}
다음에 대한 설명으로 옳지 않은 것은?
\begin{framedtext}
두 위인전 A와 B에 대한 독점 판매권을 갖고 있는 출판사가 있다. 각 위인전에 대한 소비자는 두 유형 H와 L로서 각각 50명이며, 지불용의 가격은 다음과 같다.

{\small\centering
\begin{tabular}{|c|c|c|}
\hline
\makecell{위인전 \textbackslash{} 소비자 유형} & H & L \\
\hline
A & 800원 & 400원 \\
\hline
B & 400원 & 600원 \\
\hline
\end{tabular}\par}

출판사의 한계비용은 0원이다. 출판사는 위인전 A와 B를 개별판매할지, A와 B를 함께 묶어 결합판매할지를 고려하고 있다.
\end{framedtext}
\begin{choices}
  \item 개별판매만 하는 경우 판매수입을 극대화하기 위한 가격은 A = 800원, B = 600원이다.
  \item A = 400원, B = 400원으로 가격을 책정하여 개별판매할 경우 소비자잉여가 발생한다.
  \item 결합판매만 하는 경우 두 권 묶음의 가격을 1,000원으로 책정할 때 판매수입이 극대화된다.
  \item 두 권 묶음의 가격을 1,000원으로 책정할 때 소비자잉여는 H 유형에만 발생한다.
  \item 결합판매의 최대 판매수입은 개별판매의 최대 판매수입보다 더 많다.
\end{choices}
\end{question}
\answer{1}
